% Options for packages loaded elsewhere
\PassOptionsToPackage{unicode}{hyperref}
\PassOptionsToPackage{hyphens}{url}
\PassOptionsToPackage{dvipsnames,svgnames,x11names}{xcolor}
\documentclass[
  10pt,
  fontset=fandol]{ctexart}
\usepackage{xcolor}
\usepackage[margin=16mm]{geometry}
\usepackage{amsmath,amssymb}
\setcounter{secnumdepth}{-\maxdimen} % remove section numbering
\usepackage{iftex}
\ifPDFTeX
  \usepackage[T1]{fontenc}
  \usepackage[utf8]{inputenc}
  \usepackage{textcomp} % provide euro and other symbols
\else % if luatex or xetex
  \usepackage{unicode-math} % this also loads fontspec
  \defaultfontfeatures{Scale=MatchLowercase}
  \defaultfontfeatures[\rmfamily]{Ligatures=TeX,Scale=1}
\fi
\usepackage{lmodern}
\ifPDFTeX\else
  % xetex/luatex font selection
\fi
% Use upquote if available, for straight quotes in verbatim environments
\IfFileExists{upquote.sty}{\usepackage{upquote}}{}
\IfFileExists{microtype.sty}{% use microtype if available
  \usepackage[]{microtype}
  \UseMicrotypeSet[protrusion]{basicmath} % disable protrusion for tt fonts
}{}
\makeatletter
\@ifundefined{KOMAClassName}{% if non-KOMA class
  \IfFileExists{parskip.sty}{%
    \usepackage{parskip}
  }{% else
    \setlength{\parindent}{0pt}
    \setlength{\parskip}{6pt plus 2pt minus 1pt}}
}{% if KOMA class
  \KOMAoptions{parskip=half}}
\makeatother
\setlength{\emergencystretch}{3em} % prevent overfull lines
\providecommand{\tightlist}{%
  \setlength{\itemsep}{0pt}\setlength{\parskip}{0pt}}
\usepackage{amsmath,amssymb,mathtools,needspace}
\xeCJKDeclareCharClass{CJK}{"2032,"0400->"04FF,"2160->"216B,"2460->"2473,"25A0,"25C7,"25CB,"25CF}
\IfFontExistsTF{Arial Unicode MS}{\xeCJKsetup{AutoFallBack=true}\setCJKfallbackfamilyfont{\CJKrmdefault}{Arial Unicode MS}}{}
\setcounter{secnumdepth}{0}
\setlength{\emergencystretch}{2em}
\allowdisplaybreaks
\usepackage{bookmark}
\IfFileExists{xurl.sty}{\usepackage{xurl}}{} % add URL line breaks if available
\urlstyle{same}
\hypersetup{
  pdftitle={Weierstrass 定理},
  colorlinks=true,
  linkcolor={Maroon},
  filecolor={Maroon},
  citecolor={Blue},
  urlcolor={Blue},
  pdfcreator={LaTeX via pandoc}}

\title{Weierstrass 定理}
\author{}
\date{}

\begin{document}
\maketitle

\subsubsection{3.1.2 Weierstrass 定理}\label{weierstrass-ux5b9aux7406}

用插值法或 Taylor 展开求 \(f(x) \in C[a,b]\) 的逼近多项式, 在某些点上可能没有误差, 但在整个区间 \([a,b]\) 上误差可能很大, 2.7 节给出的 Runge 现象就说明了这一点. 因此, 用 \(P_{n}(x)\) 一致逼近 \(f(x)\), 首先要解决存在性问题, 即对于 \([a,b]\) 上的连续函数 \(f(x)\), 是否存在多项式 \(P_{n}(x)\) 一致收敛于 \(f(x)\)? Weierstrass 给出了下面定理.

定理 3.1 设 \(f(x) \in C[a,b]\), 则对于任何 \(\varepsilon > 0\), 总存在一个代数多项式 \(P(x)\), 使 \[
\| f(x) - P(x) \|_{\infty} < \varepsilon
\] 在 \([a,b]\) 上一致成立.

这定理已在``数学分析''中证明过. 这里需要说明的是在许多证明方法中, Bernstein 在 1912 年给出的证明是一种构造性证明. 他根据函数整体逼近的特性造出 Bernstein 多项式 \[
\begin{cases} B_{n}(f,x) = \sum_{k=0}^{n} f\left(\frac{k}{n}\right) P_{k}(x), \\ P_{k}(x) = \binom{n}{k} x^{k}(1-x)^{n-k}, \end{cases} \tag{3.1.4}
\]

并证明了 \(\lim_{n \to \infty} B_{n}(f,x) = f(x)\) 在 \([0,1]\) 上一致成立; 若 \(f(x)\) 在 \([0,1]\) 上 \(m\) 阶导数连续, 则 \[
\lim_{n \to \infty} B_{n}^{(m)}(f,x) = f^{(m)}(x).
\] 这不但证明了定理 3.1, 而且由式 (3.1.4)给出了 \(f(x)\) 的一个逼近多项式. 它与 Lagrange 插值多项式 \[
L_{n}(x) = \sum_{k=0}^{n} f\left(x_{k}\right) l_{k}(x), \quad \sum_{k=0}^{n} l_{k}(x_{k}) = 1
\] 很相似, 对于 \(B_{n}(f,x)\), 当 \(f(x)=1\) 时也有关系式 \[
\sum_{k=0}^{n} P_{k}(x) = \sum_{k=0}^{n} \binom{n}{k} x^{k}(1-x)^{n-k} = 1. \tag{3.1.5}
\]

这只要从恒等式

\begin{quote}
校注（agent 补充，原书疑误）：原页 Lagrange 插值式后的求和明确印为 \(\sum_{k=0}^{n}l_k(x_k)=1\)，此处按原文保留。由基函数性质 \(l_k(x_k)=1\)，左端应为 \(n+1\)；通常的恒等式应为 \(\sum_{k=0}^{n}l_k(x)=1\)。这不是 OCR 误识。
\end{quote}

\href{../../source-images/pdf-060.jpeg}{查看原页}

\[
(x+y)^n = \sum_{k=0}^{n} \binom{n}{k} x^k y^{n-k} \tag{3.1.6}
\] 中令 \(y=1-x\) 就可得到. 但这里当 \(x \in [0,1]\) 时还有 \(P_k(x) \geqslant 0\), 于是 \[
\sum_{k=0}^{n} |P_k(x)| = \sum_{k=0}^{n} P_k(x) = 1
\] 是有界的, 因而只要 \(|f(x)| \leqslant \delta\) 对于任意 \(x \in [0,1]\) 成立, 则 \[
|B_n(f,x)| \leqslant \max_{0 \leqslant x \leqslant 1} |f(x)| \sum_{k=0}^{n} |P_k(x)| \leqslant \delta
\] 有界, 故 \(B_n(f,x)\) 是稳定的. 至于 Lagrange 插值多项式 \(L_n(x)\), 由于 \(\sum_{k=0}^{n} |l_k(x)|\) 无界, 因而不能保证高阶插值的稳定性与收敛性. 相比之下, 多项式 \(B_n(f,x)\) 有良好的逼近性质, 但它收敛太慢, 比三次样条逼近效果差得多, 实际中很少使用.

\subsection{本节覆盖原页的校注记录}\label{ux672cux8282ux8986ux76d6ux539fux9875ux7684ux6821ux6ce8ux8bb0ux5f55}

下列记录按原页关联，含同页相邻小节的校注；计算时同时读取上文校注和完整原页。

\begin{itemize}
\tightlist
\item
  \texttt{ch03a-E01}：原书 PDF 页 59
\end{itemize}

\end{document}
